\chapter{Mise à jour et migration}
\label{ch:upgrade}

Les mises à jour de routine sont conçues pour être sûres.
L'exception est le saut 1.0.x → 1.1.0, un breaking change unique
documenté ci-dessous.

\section{Mises à jour de routine (1.y → 1.z)}

\begin{enumerate}
  \item Lisez les notes de version pour chaque version entre votre
        version courante et la cible.
  \item Prenez une \textbf{sauvegarde de base} fraîche
        (chapitre~\ref{ch:backup}).
  \item Mettez à jour le tag d'image dans
        \texttt{docker-compose.yml} (ou re-tirez \texttt{:latest}
        si vous suivez ce tag).
  \item Lancez \texttt{docker compose up -d}. Le nouveau conteneur
        boote, exécute les nouvelles migrations contre la base
        existante, et reprend le service.
  \item Surveillez les logs pour les erreurs de migration :
        \texttt{docker compose logs mybibli | grep -i migrate}.
\end{enumerate}

Les migrations sont append-only — les changements de schéma sont
forward-compatibles. Une installation 1.4 mise à jour en 1.5 garde
toutes les lignes de 1.4 et applique les nouvelles tables /
colonnes sur place.

\begin{tip}
Pour les installations homelab, le tag Docker \texttt{:latest} est
un défaut raisonnable. Pour des montages production
(bibliothèque communautaire, association), épinglez une version
explicite pour qu'un \texttt{pull} non surveillé n'amène pas une
release sur laquelle vous n'avez pas lu les notes.
\end{tip}

\section{Sauter des versions intermédiaires}

Passer d'une installation bien plus ancienne (1.3.x par exemple)
directement à la version courante est supporté et sûr pour la
base. Le moteur de migrations applique toutes les migrations
non encore appliquées dans l'ordre des timestamps au boot,
quelle que soit la distance sautée. Les migrations de schéma
sont purement additives — aucun \texttt{DROP COLUMN}, aucun
\texttt{DROP TABLE} dans aucune release — et les quatre
backfills de données existants (le seed CSRF token de 1.1.0,
le nettoyage \texttt{manually\_edited\_fields} de 1.1.8, etc.)
sont idempotents.

Ce qui reste valable sur un saut multi-versions :

\begin{itemize}
  \item \textbf{Sauvegardez avant la mise à jour.} Il n'y a pas
        de rollback automatique. Une migration qui échoue en
        cours d'exécution laisse la base dans un état
        incohérent ; la voie de récupération est la sauvegarde
        pré-mise à jour (chapitre~\ref{ch:backup}).
  \item \textbf{Les couvertures JPG pré-1.1.4 sont perdues.}
        Jusqu'à 1.1.4, le volume \texttt{mybibli\_covers}
        n'était pas déclaré ; les couvertures écrites avant
        cette version étaient perdues à chaque upgrade de
        conteneur. Sauter directement vers 1.7 ne les récupère
        pas — re-fetch via l'action admin \textbf{Bulk re-fetch
        missing covers} (ajoutée en 1.2) ou le bouton
        \textbf{Re-fetch metadata} par titre.
  \item \textbf{Le premier boot peut prendre quelques minutes.}
        Un saut qui exécute 20+ migrations contre un gros
        catalogue est silencieux côté UX — attendez que
        \texttt{/health} réponde avant de vous inquiéter.
\end{itemize}

Le saut 1.0.x → 1.1.0 est l'unique exception qui nécessite la
procédure dédiée ci-dessous.

\section{Le saut 1.0.x → 1.1.0}\label{sec:upgrade-1.1.0}

\begin{warning}
\textbf{Ceci est un breaking change unique.}\index{mise à jour 1.1.0}
Les installations 1.0.x sont sujettes au bug documenté dans
\href{https://github.com/guycorbaz/mybibli/issues/173}{\#173} :
les migrations seed créent les identifiants \texttt{admin/admin}
et \texttt{librarian/librarian} à chaque installation neuve, y
compris en production, et contournent silencieusement l'assistant
de première mise en service.
\end{warning}

Le chemin recommandé dépend de si votre installation 1.0.x détient
des données réelles :

\subsection*{Si votre installation 1.0.x est vide (aucun titre catalogué)}

Le chemin le plus simple : effacer la base et réinstaller en 1.1.0+.

\begin{enumerate}
  \item Arrêter la pile : \texttt{docker compose down}.
  \item Supprimer le volume de base :
        \texttt{docker volume rm mybibli\_mybibli\_db\_data}
        (ou l'équivalent sur votre hôte —
        \texttt{docker volume ls} les liste).
  \item Tirer l'image 1.1.0 : \texttt{docker compose pull}.
  \item Mettre à jour le tag d'image dans
        \texttt{docker-compose.yml} si vous aviez épinglé une
        version antérieure.
  \item Démarrer : \texttt{docker compose up -d}.
  \item Ouvrir \texttt{/setup} et parcourir l'assistant.
\end{enumerate}

\subsection*{Si votre installation 1.0.x détient des données réelles}

\begin{enumerate}
  \item Prenez une sauvegarde de base (chapitre~\ref{ch:backup}).
  \item Prenez une sauvegarde covers + \texttt{.env}.
  \item Connectez-vous comme l'admin seedé et changez le mot de
        passe \emph{avant} la mise à jour. La mise à jour 1.1.0
        elle-même est sûre, mais la rotation ferme la dernière
        fenêtre pendant laquelle les identifiants par défaut
        pourraient fuiter.
  \item Mettre à jour le tag d'image dans
        \texttt{docker-compose.yml} et lancer
        \texttt{docker compose up -d}.
  \item Le boot 1.1.0 détecte l'admin existant et n'active
        \emph{pas} l'assistant — l'admin avec mot de passe tourné
        reste l'admin actif.
\end{enumerate}

\subsection*{La nouvelle variable \texttt{MYBIBLI\_SEED\_DEV\_USERS}}

À partir de 1.1.0, la migration seed est gardée par
\texttt{MYBIBLI\_SEED\_DEV\_USERS}\index{MYBIBLI\_SEED\_DEV\_USERS}.
Le défaut est \texttt{0} (pas de seed). Ne la positionnez à
\texttt{1} que pour des flux de dev local ou des tests automatisés
— jamais en production. L'assistant de première mise en service
devient le chemin canonique d'onboarding.

\section{Lire les notes de version}

Chaque release mybibli publie une page GitHub Release avec :

\begin{itemize}
  \item un paragraphe résumant le changement ;
  \item l'impact migration (aucun / additif / breaking) ;
  \item un lien vers le diff contre le tag précédent ;
  \item les PDF de ce manuel (anglais + français) en pièce jointe
        une fois le build PDF de release câblé dans la CI.
\end{itemize}

Lisez les notes avant de mettre à jour. Les breaking changes
portent une étiquette \emph{Breaking} bien visible.

\section{Nouveautés de la 1.7.6}\label{sec:whats-new-1.7.6}

La version 1.7.6 est un patch ciblé : deux corrections de fonctionnalités
visibles utilisateur et trois durcissements issus du backlog de revue de
code. \emph{Pas de migration de schéma, pas de breaking change} ;
\texttt{docker compose pull \&\& docker compose up -d} suffit.

\textbf{Visible utilisateur :}

\begin{itemize}
  \item \textbf{\#331 \string— type de média modifiable sur
        \texttt{/title/:id}.} Une BD ou un manga catalogué via son
        ISBN restait verrouillé en « livre » à vie — le formulaire
        d'édition affichait \texttt{media\_type} en lecture seule.
        Le formulaire propose maintenant un menu déroulant avec
        les six valeurs (\texttt{book}, \texttt{bd}, \texttt{cd},
        \texttt{dvd}, \texttt{magazine}, \texttt{report}) ; tout
        changement déclenche une entrée \texttt{admin\_audit}
        (\texttt{title\_media\_type\_edit}) avec le couple
        avant/après. Pratique pour corriger une BD que la BnF
        a résolue comme livre au moment du scan.
  \item \textbf{\#335 \string— téléversement manuel de couverture.}
        Le fallback de la chaîne (Open Library Covers API) rate la
        plupart des BDs francophones, des mangas et des livres
        techniques de petits éditeurs : son index est centré sur le
        marché anglophone. Cliquez \emph{Téléverser une couverture}
        sous la vignette sur \texttt{/title/:id} (rôle
        Bibliothécaire+) pour déposer un JPG / PNG / WebP / GIF
        (≤ 10\,Mio). L'image passe par exactement la même chaîne
        de décodage / redimensionnement / encodage JPEG que les
        couvertures téléchargées par les providers ; le fichier sur
        disque reste compatible. \texttt{cover\_image\_url} est
        ajouté à \texttt{manually\_edited\_fields}, ce qui exclut
        le titre de la prochaine action « Récupérer les couvertures
        manquantes » \string— votre couverture choisie est conservée.
\end{itemize}

\textbf{Durcissement (sans changement visible) :}

\begin{itemize}
  \item \textbf{\#109 \string— couverture de test des cartes
        « ajouts récents ».} Les tests de rendu de la page d'accueil
        laissaient \texttt{publication\_date} et
        \texttt{cover\_image\_url} à \texttt{None}, n'exerçant jamais
        les branches \texttt{\{\% if let Some(d) = \dots \%\}}. Un
        mauvais format de date ou une balise oubliée serait passé
        inaperçu en unitaire. Nouvelle factory
        \texttt{fake\_search\_result\_full} et nouveau test de rendu.
  \item \textbf{\#40 \string— normalisation des routes exemptées
        CSRF.} La comparaison contre la liste exemptée était littérale ;
        \texttt{POST /login/} ou \texttt{POST //login} (variantes
        que le routeur Axum dispatche identiquement vers
        \texttt{POST /login}) passaient par la validation CSRF au
        lieu d'être contournées. Nouvelle fonction
        \texttt{normalize\_exempt\_path} qui retire le slash final
        et fusionne les slashes répétés \string— les formes
        percent-encodées ne sont volontairement \emph{pas} décodées
        (la tentative de bypass \texttt{/login\%2F..\%2Fadmin} reste
        rejetée).
  \item \textbf{\#36 \string— bypass du résolveur de session sur
        les routes statiques.} \texttt{session\_resolve\_middleware}
        traversait le résolveur sur \emph{chaque} requête, y compris
        \texttt{/static/*}, \texttt{/covers/*} et \texttt{/logo/*}.
        Un seul chargement de la page d'accueil déclenchait des
        dizaines de requêtes d'assets, chacune lisant une ligne
        \texttt{sessions} (et insérant une ligne anonyme sur un
        onglet vierge). Le middleware court-circuite désormais ces
        préfixes en plus du bypass \texttt{/health} existant.
\end{itemize}

\section{Nouveautés de la 1.7.5}\label{sec:whats-new-1.7.5}

La version 1.7.5 est un patch ciblé qui groupe un correctif visible
utilisateur et sept commits de durcissement issus de
code-review-findings. \emph{Aucune migration de schéma, aucune
cassure} ; \texttt{docker compose pull \&\& docker compose up -d}
habituel.

\textbf{Visible utilisateur :}

\begin{itemize}
  \item \textbf{\#325 \string— ajout / retrait de contributeur
        silencieusement cassé sur \texttt{/title/:id}.} Cliquer
        sur \emph{+ Ajouter un contributeur} sur la page de
        détail d'un titre ouvrait bien le formulaire et acceptait
        le nom et le rôle, mais \emph{Enregistrer} ne produisait
        aucun effet visible. Cause racine : le formulaire pointe
        \texttt{hx-target="\#feedback-list"}, mais le slot était
        déclaré uniquement sur \texttt{/catalog} et
        \texttt{/admin} — pas sur \texttt{/title/:id}. Le même
        défaut affectait silencieusement le bouton
        \emph{Retirer}. Régression du fix v1.7.2 qui avait ajouté
        le flux de détail sans porter le slot. Un garde-fou
        d'audit templates panique désormais si une page référence
        \texttt{hx-target="\#feedback-list"} sans déclarer
        \texttt{id="feedback-list"}.
\end{itemize}

\textbf{Durcissement (aucune cassure utilisateur) :}

\begin{itemize}
  \item \textbf{\#91 \string— re-rendu du formulaire sur erreur
        de validation.} Admin → Système → Prêts/Langue
        retournent 400 avec le formulaire re-rendu (la valeur
        saisie est préservée) et un retour d'erreur inline, au
        lieu d'un 400 vide silencieusement avalé par HTMX 2.0.
  \item \textbf{\#42 \string— jeton CSRF en premier-enfant
        strict.} L'audit \texttt{forms\_include\_csrf\_token}
        exige désormais \texttt{\_csrf\_token} comme tout premier
        \texttt{<input>} de chaque formulaire POST, pas
        seulement parmi les cinq premiers.
  \item \textbf{\#48 \string— audit des littéraux HTML Rust.}
        Nouveau test frère
        \texttt{rust\_emitted\_post\_forms\_include\_csrf\_token}
        parcourt \texttt{src/**/*.rs} et refuse tout
        \texttt{<form method="POST">} en chaîne littérale qui ne
        déclare pas son jeton CSRF à proximité. Verrouille le
        site existant \texttt{src/routes/locations.rs:301} ;
        tout nouveau formulaire émis depuis Rust devra déclarer
        son jeton.
  \item \textbf{\#220 \string— HX-Trigger en forme JSON.}
        \texttt{modal\_confirm\_retarget\_guard} tolère désormais
        la forme JSON-objet du header \texttt{HX-Trigger} en
        plus de la forme virgule-séparée. Défensif : l'émetteur
        actuel utilise la forme virgule.
  \item \textbf{\#31 \string— scanner-guard respect du pipeline
        d'entrée.} Le JavaScript scanner-guard qui transmet les
        rafales de scanner USB vers le champ texte d'un modal
        ouvert respecte désormais \texttt{maxLength} et passe
        outre pendant une composition IME active. Les sous-items
        3/4/5 étaient déjà classés N/A à la forme actuelle des
        modaux.
  \item \textbf{\#90 \string— DRY admin\_system.} La requête
        UPDATE de verrou optimiste dupliquée entre
        \texttt{save\_setting} et \texttt{run\_provider\_update}
        est consolidée ; le second délègue désormais et
        map\_errs le message de conflit pour interpoler le label
        du provider.
  \item \textbf{\#152 \string— test \texttt{/series} anonyme.}
        Ajoute le test de garde de rôle manquant qui verrouille
        le contrat FR95 d'accès anonyme à \texttt{/series}
        (miroir de \texttt{/catalog} et \texttt{/locations}).
\end{itemize}

\section{Rollback}

Le rollback \emph{n'est pas} un flux de première classe. Les
migrations de schéma avancent uniquement — il n'y a pas de
\texttt{down.sql}. La stratégie de récupération supportée est :

\begin{enumerate}
  \item Restaurer la sauvegarde de base pré-mise-à-jour
        (chapitre~\ref{ch:backup}).
  \item Réépingler le tag d'image à la version antérieure.
  \item Démarrer la pile.
\end{enumerate}

Cela fonctionne car la sauvegarde est cohérente avec le schéma
antérieur. Tenter de pointer une image plus ancienne sur une base
fraîchement migrée échoue au boot.
